\section{Capítulo~\ref{sec:gaba}. \nt{GLUT} e \nt{GABA} juntos}

As afirmações lógicas e dinâmicas do capítulo são deduzidas nele mesmo, por análise linear e pela lei de Ohm; a experiência mostra que os circuitos em que elas se apoiam (veto com atraso, inibição direta logo após a excitação, estabilização pela inibição, ritmo do laço \nt{GLUT}--\nt{GABA}) de fato funcionam no cérebro.

{\footnotesize
\setlength{\tabcolsep}{3pt}\renewcommand{\arraystretch}{1.15}
\begin{longtable}{>{\raggedright}p{0.5cm}>{\raggedright}p{4.0cm}>{\raggedright}p{5.6cm}>{\raggedright}p{2.3cm}>{\raggedright\arraybackslash}p{1.7cm}}
\toprule
N.º & Proposição & Experimento e o que foi medido & Fonte & Status \\
\midrule\endfirsthead
\toprule
N.º & Proposição & Experimento e o que foi medido & Fonte & Status \\
\midrule\endhead
1 & Os neurônios \nt{GABA} são de um quinto a um quarto dos neurônios do córtex & Imunomarcação para GABA em $10$ áreas do córtex de macaco: cerca de $25\,\%$ dos neurônios em $8$ áreas, menos de $20\,\%$ no córtex visual primário. Revisão da diversidade dos neurônios inibitórios do córtex & \cite{Hendry1987,Markram2004} & medido \\
2 & O que a inibição faz é decidido em grande parte pelo local de contato: dendrito, corpo, local de nascimento do disparo, terminal do fio de outro neurônio & Revisão: as classes de neurônios \nt{GABA} do córtex diferem pelo alvo no destinatário. Registro simultâneo do corpo e do dendrito de um neurônio piramidal: a inibição direta é muito mais forte no corpo do que nos dendritos. A inibição no terminal do fio de outro neurônio, na medula espinal, reduz a liberação de neurotransmissor de uma única linha de entrada (revisão das medições) & \cite{Markram2004,Pouille2001,Rudomin1999} & medido \\
3 & A troca de sinal por um intermediário custa um atraso, cerca de um milissegundo; a inibição que chega logo após a excitação estreita a janela de coincidência & Neurônios piramidais do hipocampo de rato em fatias: a inibição por um intermediário chega logo após a excitação direta com atraso curto, e a janela em que duas entradas excitatórias se somam até o limiar é menor que $2\ms$. Nos dendritos, onde essa inibição é mais fraca, a janela é mais larga & \cite{Pouille2001} & medido \\
4 & O veto $a\wedge\bar b$~\eqref{eq:veto} com OU e a constante um é funcionalmente completo (proposição~\ref{prop:gabacomplete}) & Demonstração no capítulo. Resultado clássico: uma rede de elementos de limiar com entradas inibitórias realiza qualquer expressão lógica. Afirmação sobre o modelo, sem experimento & \cite{McCullochPitts1943} & teoria \\
5 & O veto com atraso dá um detector de direção do movimento com velocidade ótima $v^*=\Delta x/d$~\eqref{eq:vstar} & Retina do coelho: a seletividade à direção é criada por uma inibição que, no movimento na direção ``nula'', anula a excitação. O registro separado das entradas excitatória e inibitória das células seletivas mostrou que as células amácrinas estreladas (\nt{GABA}) as inibem de forma assimétrica, só pelos prolongamentos voltados para o lado ``nulo''. A fórmula de $v^*$ é dedução do capítulo & \cite{Barlow1965,Fried2002} & medido \\
6 & A inibição por derivação divide a tensão em vez de subtrair~\eqref{eq:shunt} (fig.~\ref{fig:shunt}) & Lei de Ohm para um divisor de três condutâncias, com o potencial de equilíbrio do cloreto perto do repouso; para um neurônio sem disparos & dedução no capítulo & teoria \\
7 & Sobre a frequência de disparos a derivação age por subtração, e a divisão vem da inibição junto com ruído de fundo equilibrado & Modelo: num neurônio que dispara, a corrente pela derivação é quase constante, e o efeito sobre a frequência é subtrativo. Experimento: em neurônios piramidais do córtex somatossensorial de rato (fatias) introduziram-se, por fixação dinâmica, condutâncias de fundo excitatórias e inibitórias; o aumento simultâneo e equilibrado de ambas divide a inclinação da resposta sem deslocamento notável da frequência. Revisão: a divisão pela atividade total ocorre em muitas regiões do cérebro & \cite{HoltKoch1997,Chance2002,Carandini2012} & medido \\
8 & Com $\beta>1$ a inibição mútua escolhe um único vencedor (proposição~\ref{prop:wta}, \eqref{eq:wta}) & Os autovalores $-(1\pm\beta)/\tau$ são dedução do capítulo. As frequências $0.73$ e $0.53$ com $\beta=0.5$ são o ponto fixo $r_{1,2}=(I_{1,2}-\beta I_{2,1})/(1-\beta^2)$; não há script em \texttt{models/}. O capítulo não cita experimento direto & dedução no capítulo & teoria \\
9 & Reset da memória por sinal via \nt{GABA}; dois grupos com inibição mútua formam um latch & Decorre da fig.~\ref{fig:bistable} e da proposição~\ref{prop:wta}. Sem experimento & dedução no capítulo & teoria \\
10 & O equilíbrio do laço \nt{GLUT}--\nt{GABA} é estável se a inibição for forte e rápida o bastante (proposição~\ref{prop:ei}) & Modelo de duas populações~\eqref{eq:ei}; as condições (i) e (ii) são o determinante e o traço de sua matriz, deduzidas no capítulo & \cite{WilsonCowan1972} & teoria \\
11 & Com $w_{EE}=1.5$ e $w_{EI}w_{IE}=2$, a inibição rápida ($\tau_I=2\ms$) mantém $E^*=0.67h$, e a lenta ($30\ms$) provoca oscilações (fig.~\ref{fig:ei}) & Solução numérica de~\eqref{eq:ei}, script \texttt{figures/make\_ei.py} & cálculo & modelo \\
12 & O córtex trabalha no regime de estabilização pela inibição; a assinatura: empurre os neurônios \nt{GABA}, e sua frequência cai~\eqref{eq:isn} & A teoria previu a resposta paradoxal. Experimento: registros intracelulares no córtex visual primário do gato; quando a resposta é suprimida por um estímulo ao redor, a condutância inibitória não cresce, mas cai junto com a excitatória. A excitação optogenética de neurônios PV do camundongo nos córtices visual, somatossensorial e motor reduz a frequência dos neurônios inibitórios, inclusive sem estímulo e sob anestesia leve. O coeficiente $-1/3$ com os números da fig.~\ref{fig:ei} é dedução do capítulo & \cite{Tsodyks1997isn,Ozeki2009,Sanzeni2020} & medido \\
13 & O laço ``\nt{GLUT} excita \nt{GABA}, \nt{GABA} inibe \nt{GLUT}'' gera um ritmo rápido com período de $12$--$50\ms$ & A estimulação optogenética de neurônios \nt{GABA} rápidos no córtex somatossensorial do camundongo a $8$--$200$\,Hz amplifica seletivamente o ritmo gama ($20$--$80$\,Hz, período de $12$--$50\ms$); a estimulação de neurônios \nt{GLUT} amplifica só ritmos mais lentos & \cite{Cardin2009,Buzsaki2012} & medido \\
\bottomrule
\end{longtable}}
